\chapter{समस्याएँ हल करना}\label{ch:g3:problems}

चारों संक्रियाएँ जानना एक बात है; यह जानना कि कहानी को
\emph{कौन-सी} चाहिए, दूसरी बात। यह अध्याय शब्द-समस्याओं के लिए
एक तरीक़ा देता है --- हर बार वही चार कदम --- और सारणियों तथा
चित्रालेखों से जानकारी पढ़ना सिखाता है।

\section{चार कदम}

\begin{method}[शब्द-समस्या हल करना]\label{met:g3:problems:steps}
\begin{enumerate}
  \item समस्या दो बार \emph{पढ़ो}; अपने शब्दों में कहो कि क्या
        पूछा गया है और क्या पता है।
  \item संक्रिया \emph{चुनो}: चुनाव साफ़ न हो तो झट से एक
        रेखाचित्र या पट्टी-चित्र बना लो।
  \item \emph{निकालो}, खानों में या मन-ही-मन, और सुरक्षा के लिए
        एक अंदाज़ा साथ रखो।
  \item पूरे वाक्य में \emph{उत्तर} दो, मात्रक के साथ --- और
        देखो कि उत्तर तुक का है (कोई बच्चा
        $300$ kg का नहीं होता)।
\end{enumerate}
\end{method}

\begin{example}[कौन-सी संक्रिया?]\label{ex:g3:problems:which}
शब्द संक्रिया का इशारा देते हैं, रेखाचित्र फ़ैसला करता है:
\begin{itemize}
  \item \emph{मिलाना / कुल कितने}: जोड़;
  \item \emph{निकालना / कितने ज़्यादा / कितना कम है}:
        घटाव;
  \item \emph{इतने-इतने के इतने समूह / इतनी बार}:
        गुणा;
  \item \emph{बराबर बाँटना / कितने समूह}: भाग।
\end{itemize}
\end{example}

\begin{omfigure}
\begin{tikzpicture}[scale=0.62]
  % जोड़ की पट्टी
  \draw[thick, fill=omDef!20] (0,0) rectangle (3.4,0.8);
  \draw[thick, fill=omThm!20] (3.4,0) rectangle (5.8,0.8);
  \node[font=\small] at (1.7,0.4) {$26$};
  \node[font=\small] at (4.6,0.4) {$17$};
  \draw[Stealth-Stealth] (0,-0.35) -- (5.8,-0.35);
  \node[below, font=\small] at (2.9,-0.4) {? कुल: $26 + 17$};
  % तुलना की पट्टी
  \begin{scope}[xshift=9.4cm]
  \draw[thick, fill=omDef!20] (0,0.9) rectangle (5.2,1.7);
  \node[font=\small] at (2.6,1.3) {$52$};
  \draw[thick, fill=omThm!20] (0,0) rectangle (3.5,0.8);
  \node[font=\small] at (1.75,0.4) {$35$};
  \draw[Stealth-Stealth] (3.5,0.4) -- (5.2,0.4);
  \node[below, font=\small] at (2.6,-0.15)
    {? ज़्यादा: $52 - 35$};
  \end{scope}
\end{tikzpicture}

{\small पट्टी-चित्र संक्रिया को दिखा देते हैं: अगल-बगल रखे
हिस्से जोड़ माँगते हैं; दो पट्टियों की तुलना घटाव माँगती है।}
\end{omfigure}

\begin{example}[दो कदमों वाली समस्या]\label{ex:g3:problems:twostep}
``माया $6$-$6$ बोतलों के $3$ पैकेट ख़रीदती है और $2$ बोतलें पी
जाती है। कितनी बोतलें बचीं?''
\begin{enumerate}
  \item क्या पूछा है: बची हुई बोतलों की संख्या। पता है: $6$-$6$
        के $3$ पैकेट; $2$ पी गईं।
  \item रेखाचित्र: $6$-$6$ के $3$ समूह, फिर $2$ निकालो। दो
        संक्रियाएँ!
  \item निकालो: $3 \times 6 = 18$, फिर $18 - 2 = 16$।
  \item वाक्य: $16$ बोतलें बचीं।
\end{enumerate}
पहली संक्रिया के बाद कभी मत रुको: सवाल दोबारा पढ़कर देखो कि
कहानी पूरी हुई या नहीं।
\end{example}

\begin{example}[एक बेकार संख्या]\label{ex:g3:problems:trap}
``एक नानबाई, जो $52$ साल का है, $120$ बन पकाता है और $85$ बेच
देता है। कितने बन बचे?'' उम्र सिर्फ़ ध्यान भटकाने के लिए है:
उत्तर के लिए बस $120 - 85 = 35$ बन चाहिए। गणना से पहले हमेशा
संख्याएँ छाँटो: सवाल सचमुच किनका उपयोग करता है?
\end{example}

\section{सारणियाँ और चित्रालेख पढ़ना}

\begin{example}[एक सारणी]\label{ex:g3:problems:table}
पुस्तकालय से ली गई किताबें:
\begin{center}
\renewcommand{\arraystretch}{1.3}
\begin{tabular}{l|cccc}
 & चित्रकथा & उपन्यास & विज्ञान & कुल \\
\hline
सप्ताह 1 & $14$ & $9$ & $6$ & $29$ \\
सप्ताह 2 & $11$ & $12$ & $8$ & $31$ \\
\end{tabular}
\end{center}
पढ़ना: सप्ताह 2, उपन्यास: $12$। सारणी से गणना: दोनों सप्ताहों
में मिलाकर $14 + 11 = 25$ चित्रकथाएँ ली गईं; सबसे व्यस्त
सप्ताह दूसरा था ($31 > 29$)।
\end{example}

\begin{example}[एक चित्रालेख]\label{ex:g3:problems:pictogram}
चित्रालेख में एक चिह्न एक तय मात्रा दिखाता है --- पहले कुंजी
पढ़ो!

\begin{omfigure}
\begin{tikzpicture}[scale=0.62]
  \node[left, font=\small] at (0,2.4) {सोमवार};
  \foreach \i in {0,1,2} \fill[omDef] (0.5+\i*0.8,2.4) circle (0.26);
  \node[left, font=\small] at (0,1.2) {मंगलवार};
  \foreach \i in {0,1,2,3,4} \fill[omDef] (0.5+\i*0.8,1.2) circle (0.26);
  \node[left, font=\small] at (0,0) {बुधवार};
  \foreach \i in {0,1} \fill[omDef] (0.5+\i*0.8,0) circle (0.26);
  \fill[omDef] (0.5+1.6,0) circle (0.26);
  \node[right, font=\small] at (5.4,1.2)
    {कुंजी: एक बिंदु $= 4$ केक बिके};
\end{tikzpicture}

{\small एक चित्रालेख: सोमवार को $3$ बिंदु हैं, यानी
$3 \times 4 = 12$ केक; मंगलवार $5$ बिंदु, $20$ केक; बुधवार
$3$ बिंदु, $12$ केक।}
\end{omfigure}
\end{example}

\section{अभ्यास}

\begin{exercise}[$\star$]\label{exo:g3:problems:1}
हर कहानी के लिए केवल संक्रिया का \emph{नाम} बताओ (गणना मत करो):
$48$ भेड़ों का झुंड $37$ भेड़ों के झुंड से मिल जाता है; \ $56$
चेरी $7$ बच्चों में बाँटी गईं; \ $8$-$8$ पेंसिलों के $9$ डिब्बे;
\ अली के पास $71$ कार्ड थे और उसने $18$ दे दिए।
\end{exercise}

\begin{exercise}[$\star$]\label{exo:g3:problems:2}
चार कदमों से हल करो: ``एक पार्किंग में सुबह $246$ कारें हैं;
दोपहर को $158$ और आ जाती हैं। अब कितनी कारें हैं?''
\end{exercise}

\begin{exercise}[$\star$]\label{exo:g3:problems:3}
हल करो: ``$90$ m की रस्सी $9$-$9$ m के टुकड़ों में काटी जाती
है। कितने टुकड़े?''
\end{exercise}

\begin{exercise}[$\star$]\label{exo:g3:problems:4}
हल करो: ``एक टिकट $8$ का है। $26$ टिकटों के कितने लगेंगे?''
\end{exercise}

\begin{exercise}[$\star$]\label{exo:g3:problems:5}
पट्टी-चित्र बनाओ, फिर हल करो: ``एम्मा के पास $64$ स्टिकर हैं,
ह्यूगो के पास $29$ कम। ह्यूगो के पास कितने हैं?''
\end{exercise}

\begin{exercise}[$\star$]\label{exo:g3:problems:6}
दो कदम: ``एक पेटी में $45$ सेब आते हैं। एक दुकानदार को $6$
पेटियाँ मिलती हैं और वह $178$ सेब बेच देता है। कितने सेब बचे?''
\end{exercise}

\begin{exercise}[$\star$]\label{exo:g3:problems:7}
\cref{ex:g3:problems:table} की सारणी से: दोनों सप्ताहों में
मिलाकर विज्ञान की कितनी किताबें ली गईं? सप्ताह 2 में सप्ताह 1
से कितने ज़्यादा उपन्यास?
\end{exercise}

\begin{exercise}[$\star$]\label{exo:g3:problems:8}
\cref{ex:g3:problems:pictogram} के चित्रालेख से: तीनों दिनों में
मिलाकर कितने केक बिके? किस दिन सबसे ज़्यादा केक
बिके?
\end{exercise}

\begin{exercise}[$\star$]\label{exo:g3:problems:9}
इस समस्या में एक बेकार संख्या है --- उसे ढूँढ़ो, फिर हल करो:
``$210$ पन्नों की किताब $12$ की है। लियो रोज़ $35$ पन्ने पढ़ता
है। किताब पूरी करने में उसे कितने दिन लगेंगे?''
\end{exercise}

\begin{exercise}[$\star\star$]\label{exo:g3:problems:10}
``$2$ बड़ों और $3$ बच्चों का एक परिवार संग्रहालय देखने जाता है।
बड़ों का टिकट $9$ का है, बच्चों का $4$ का। वे $50$ का नोट देते
हैं। उन्हें कितने पैसे वापस मिलते हैं?'' (तीन कदम: बड़े, बच्चे,
फिर बाक़ी पैसे।)
\end{exercise}

\begin{exercise}[$\star\star$]\label{exo:g3:problems:11}
दो कदमों वाली अपनी एक समस्या बनाओ जिसका उत्तर $100$ हो, चार
कदमों के साथ उसका हल लिखो, और किसी सहपाठी पर आज़माओ।
\end{exercise}
